\documentclass[UTF8,a4paper,zihao=-4]{ctexart}
\usepackage[margin=2.3cm]{geometry}
\setmainfont{Times New Roman}
\setsansfont{Arial}
\setmonofont{Consolas}
\usepackage{amsmath,amssymb,graphicx,booktabs,longtable,array,calc,multirow}
\usepackage{hyperref,xurl,float,caption,fancyvrb,seqsplit}
\hypersetup{hidelinks}
\usepackage{fvextra}
\DefineVerbatimEnvironment{Highlighting}{Verbatim}{breaklines,commandchars=\\\{\}}
\newcommand{\tightlist}{\setlength{\itemsep}{0pt}\setlength{\parskip}{0pt}}
\setcounter{tocdepth}{1}
\setlength{\emergencystretch}{4em}
\setlength{\parindent}{2em}
\setlength{\parskip}{0.25em}
\renewcommand{\arraystretch}{1.15}
\captionsetup{font=small}
\pagestyle{plain}
\makeatletter
\def\maxwidth{\ifdim\Gin@nat@width>\linewidth\linewidth\else\Gin@nat@width\fi}
\def\maxheight{\ifdim\Gin@nat@height>0.78\textheight0.78\textheight\else\Gin@nat@height\fi}
\makeatother
\setkeys{Gin}{width=\maxwidth,height=\maxheight,keepaspectratio}
\begin{document}
\pagenumbering{gobble}
\begin{center}{\zihao{3}\bfseries 观测语义约束下的啤酒风味预测、化合物筛选与条件生成}\end{center}
\begin{center}\textbf{摘要}\end{center}
针对啤酒化学数据与感官评分的尺度混合、Top-3共线冗余及小样本风格迁移问题，构建由观测语义、训练内选择和独立风险验证衔接的三问方案。将2197条原始品评记录映射到250款啤酒，定义评分员等权0—6面板目标，并保留评分卡层次和重复品评敏感性。评分卡交付为250行×53列、无缺失值：1列啤酒编号，50列由冻结模型预测的感官评分，另2列acetaldehyde为阈值锚定推断值（无面板0—6记录，由附件1实测浓度经log10反变换后与附件8的50.0mg/L感官检出阈值比较，250款酒全部低于阈值故取0=Not present），后两列明确标注为非面板评分。该评分卡存在一项已量化并披露的局限：预测等级全部落在0—4级，真值≥5的121个单元被系统低估（平均低估2.07，连续预测上限仅4.34），故整体MAE可用而高强度尾部不可用。冻结的Extra Trees取得MAE 0.4454、RMSE 0.6305，层次网络未超过强传统基线；非香气强度项的平均MAE（0.3486）优于香气强度项（0.5047）。化合物筛选在训练内部比较Spearman、2+1、三步条件前向与冗余控制，最终选出2+1并统一50×3结果；包含方法选择的外层MAE为0.5024，揭示小样本选法的方差代价。对四种既有生成模型完成22风格分布、分类一致性、相关结构及联合低端模式诊断：交付生成器沿用诊断之前冻结的验证折规则而非条件SWD排名，因为CVAE的表面SWD优势与测试集大小强相关（r=0.70，仅在1—2款测试酒的风格上领先，5款以上时优势消失）。9款独立源酒的四策略配对比较按两套缺陷口径报告：在赛题所指的超阈值口径上，硬约束把越阈次数由无约束的7次降到0次；在任意微量上升的严格口径上仍有4—5/9款酒上升，该残余风险列为未解决局限。结果支持保留拒绝选项、报告个体风险并区分编码约束与真实化学保证。三组封闭实验完成后停止模型探索，交付可复现代码、风格谱、候选因素和完整验证材料。本文为如实披露AI贡献的研究定稿，队员独立核心工作状态另列。

\noindent\textbf{关键词：}面板评分；训练内选择；替代组稳定性；条件生成；独立风险

\vfill\noindent\small AI辅助研究参考稿。核心建模与分析尚待参赛队独立完成，不作为独立原创正式提交稿。
\newpage\normalsize\tableofcontents\newpage
\pagenumbering{arabic}
\input{build_rebuild/body.tex}
\clearpage\appendix
\input{build_rebuild/appendix.tex}
\section{复现环境与全部新增程序}
以下为本轮新增计算程序，均含AI辅助说明。原始数据未修改。运行命令：Q2/Q3 封闭实验与论文重生成见 scripts/run\_final.ps1（scripts/run\_rebuild.ps1 为其兼容转发入口）；Q1 从源码完整重建见 scripts/run\_q1.ps1。脚本依赖见 configs/rebuild\_requirements.txt。

\subsection{audit\_extended.py}
\VerbatimInput[fontsize=\tiny,breaklines,breakanywhere]{../src/rebuild/audit_extended.py}

\subsection{closed.py}
\VerbatimInput[fontsize=\tiny,breaklines,breakanywhere]{../src/rebuild/closed.py}

\subsection{closed\_figures.py}
\VerbatimInput[fontsize=\tiny,breaklines,breakanywhere]{../src/rebuild/closed_figures.py}

\subsection{closed\_report.py}
\VerbatimInput[fontsize=\tiny,breaklines,breakanywhere]{../src/rebuild/closed_report.py}

\subsection{closed\_tables.py}
\VerbatimInput[fontsize=\tiny,breaklines,breakanywhere]{../src/rebuild/closed_tables.py}

\subsection{closed\_verify.py}
\VerbatimInput[fontsize=\tiny,breaklines,breakanywhere]{../src/rebuild/closed_verify.py}

\subsection{figures.py}
\VerbatimInput[fontsize=\tiny,breaklines,breakanywhere]{../src/rebuild/figures.py}

\subsection{paths.py}
\VerbatimInput[fontsize=\tiny,breaklines,breakanywhere]{../src/rebuild/paths.py}

\subsection{q1\_acetaldehyde.py}
\VerbatimInput[fontsize=\tiny,breaklines,breakanywhere]{../src/rebuild/q1_acetaldehyde.py}

\subsection{q1\_complete.py}
\VerbatimInput[fontsize=\tiny,breaklines,breakanywhere]{../src/rebuild/q1_complete.py}

\subsection{q1\_grade\_range.py}
\VerbatimInput[fontsize=\tiny,breaklines,breakanywhere]{../src/rebuild/q1_grade_range.py}

\subsection{q1\_scale\_check.py}
\VerbatimInput[fontsize=\tiny,breaklines,breakanywhere]{../src/rebuild/q1_scale_check.py}

\subsection{q2.py}
\VerbatimInput[fontsize=\tiny,breaklines,breakanywhere]{../src/rebuild/q2.py}

\subsection{q3.py}
\VerbatimInput[fontsize=\tiny,breaklines,breakanywhere]{../src/rebuild/q3.py}

\subsection{q3\_diagnostics.py}
\VerbatimInput[fontsize=\tiny,breaklines,breakanywhere]{../src/rebuild/q3_diagnostics.py}

\subsection{report.py}
\VerbatimInput[fontsize=\tiny,breaklines,breakanywhere]{../src/rebuild/report.py}

\subsection{run.py}
\VerbatimInput[fontsize=\tiny,breaklines,breakanywhere]{../src/rebuild/run.py}

\subsection{supplement.py}
\VerbatimInput[fontsize=\tiny,breaklines,breakanywhere]{../src/rebuild/supplement.py}

\subsection{verify.py}
\VerbatimInput[fontsize=\tiny,breaklines,breakanywhere]{../src/rebuild/verify.py}

\subsection{run\_final.ps1}
\VerbatimInput[fontsize=\tiny,breaklines,breakanywhere]{../scripts/run_final.ps1}

\subsection{run\_q1.ps1}
\VerbatimInput[fontsize=\tiny,breaklines,breakanywhere]{../scripts/run_q1.ps1}

\subsection{run\_rebuild.ps1}
\VerbatimInput[fontsize=\tiny,breaklines,breakanywhere]{../scripts/run_rebuild.ps1}

\subsection{build\_rebuild\_pdf.py}
\VerbatimInput[fontsize=\tiny,breaklines,breakanywhere]{../scripts/build_rebuild_pdf.py}

\subsection{final\_quality.py}
\VerbatimInput[fontsize=\tiny,breaklines,breakanywhere]{../scripts/final_quality.py}

\subsection{package\_final.py}
\VerbatimInput[fontsize=\tiny,breaklines,breakanywhere]{../scripts/package_final.py}

\end{document}